% ============================================================
% Title:  UMG TikZ lexicon: visual primitives of the Unified
%         Model Graph (UMG) system
% File:   umg-style.tex
% Author: H.-T. Yu
% Date:   2026-06-12
% ============================================================
% Input this file in the preamble of any document (manuscript or
% standalone figure) that draws UMG diagrams. Requires TikZ.
%
% Channel-to-attribute mapping (see Table 2 of the manuscript):
%   Fill (luminance) -> observability  (shaded = observed; white = latent)
%   Shape            -> support        (circle = continuous; square = categorical)
%   Border           -> role           (single = random variable; double = deterministic)
%   Diamond          -> fixed unknown parameter (frequentist target of estimation)
%   Triangle         -> known constant (e.g., the unit constant for means)
%   Rounded box      -> plate (exchangeable replication; nesting = hierarchy)
% ------------------------------------------------------------
\usetikzlibrary{arrows.meta,positioning,fit,backgrounds,shapes.geometric,calc}

\tikzset{
  % ----- Vertex primitives -----------------------------------------
  % continuous latent random variable: white circle
  vCL/.style={draw, circle, minimum size=9.5mm, inner sep=1pt, line width=0.7pt, fill=white},
  % continuous observed random variable: shaded circle
  vCO/.style={draw, circle, minimum size=9.5mm, inner sep=1pt, line width=0.7pt, fill=black!13},
  % categorical latent random variable (e.g., latent class): white square
  vDL/.style={draw, rectangle, minimum size=8.5mm, inner sep=2pt, line width=0.7pt, fill=white},
  % categorical observed random variable (e.g., item response): shaded square
  vDO/.style={draw, rectangle, minimum size=8.5mm, inner sep=2pt, line width=0.7pt, fill=black!13},
  % fixed unknown parameter: open diamond
  vP/.style={draw, diamond, aspect=1.15, minimum size=7.5mm, inner sep=0.5pt, line width=0.7pt, fill=white},
  % deterministic node: double border (shape follows support)
  vDetC/.style={draw, circle, double, double distance=1.1pt, minimum size=9.5mm, inner sep=1pt, line width=0.5pt, fill=white},
  vDetD/.style={draw, rectangle, double, double distance=1.1pt, minimum size=8.5mm, inner sep=2pt, line width=0.5pt, fill=white},
  % known constant: triangle (RAM heritage); label inside
  vK/.style={draw, isosceles triangle, isosceles triangle apex angle=60,
             shape border rotate=90, minimum size=7mm, inner sep=0.5pt,
             line width=0.7pt, fill=white},
  % small disturbance/residual node (continuous latent, reduced size)
  vEps/.style={draw, circle, minimum size=6.5mm, inner sep=0.5pt, line width=0.6pt, fill=white},
  % ----- Edge primitives -------------------------------------------
  % stochastic/functional dependence (regression, loading)
  eReg/.style={->, >={Stealth[length=2.8mm]}, line width=0.7pt},
  % covariance / unanalyzed association (symmetric)
  eCov/.style={<->, >={Stealth[length=2.5mm]}, line width=0.7pt},
  % deterministic assignment
  eDet/.style={->, >={Stealth[length=2.8mm]}, line width=0.7pt, dashed},
  % mixing edge: latent class -> parameter set (open arrowhead)
  eMix/.style={->, >={Triangle[open, length=3mm]}, line width=0.7pt},
  % variance stub: double-headed self-loop on a node
  eVar/.style={<->, >={Stealth[length=2mm]}, line width=0.6pt, looseness=5},
  % edge parameter label
  elab/.style={font=\small, fill=white, inner sep=1pt, sloped=false},
  elabplain/.style={font=\small, inner sep=1.5pt},
  % ----- Plates ------------------------------------------------------
  plate/.style={draw, rounded corners=4pt, inner sep=9pt, line width=0.6pt},
  platelab/.style={anchor=south east, inner sep=4pt, font=\small\itshape},
  platelabnw/.style={anchor=north west, inner sep=4pt, font=\small\itshape},
  % ----- Auxiliary ---------------------------------------------------
  dist/.style={font=\footnotesize, text=black!70, inner sep=1pt},
  causalbadge/.style={draw, rounded corners=2pt, font=\footnotesize\sffamily,
                      inner sep=2.5pt, fill=black!5},
}

% Residual stub: arrow into #1 from direction #2 (angle), with label #3.
% Compact abbreviation for an explicit disturbance node; see Section 4.5.
\newcommand{\umgstub}[3]{%
  \draw[eReg] ($(#1)+(#2:0.95)$) -- (#1);
  \node[elabplain] at ($(#1)+(#2:1.25)$) {#3};%
}
